\section{玻色--爱因斯坦凝聚与超流}

\subsection{理想玻色气的 BEC}

三维理想玻色气在
\begin{equation}
  T_c
  =\frac{2\pi\hbar^2}{mk_B}
  \Bigl(\frac{n}{\zeta(3/2)}\Bigr)^{2/3}
\end{equation}
发生动量 $\mathbf{k}=0$ 态的宏观占据：$n_0/n>0$。这是 $U(1)$ 相位对称性破缺的最简实现（严格说理想气需小心处理）。

\subsection{弱相互作用：Bogoliubov 理论}

接触相互作用 $g$ 下，凝聚体密度 $\rho_c$，对涨落做 Bogoliubov 变换
\begin{equation}
  \begin{pmatrix}\psi_{\mathbf{k}}\\\psi_{-\mathbf{k}}^\dagger\end{pmatrix}
  =
  \begin{pmatrix}u_k&v_k^*\\v_k&u_k^*\end{pmatrix}
  \begin{pmatrix}\varphi_{\mathbf{k}}\\\varphi_{-\mathbf{k}}^\dagger\end{pmatrix},
  \qquad
  |u|^2-|v|^2=1,
\end{equation}
对角化后得到色散
\begin{equation}
  E_k
  =\sqrt{
  \varepsilon_k(\varepsilon_k+2g\rho_c)
  },
  \qquad
  \varepsilon_k=\frac{\hbar^2k^2}{2m}.
\end{equation}
长波 $E_k\approx \hbar c |k|$，$c=\sqrt{g\rho_c/m}$ 为声速——这是超流临界速度的微观起源之一（Landau 判据）。

\subsection{超流性}

超流密度 $\rho_s$ 刻画对相位梯度的刚性响应。二维有限温下由 BKT 机制通过涡旋解离破坏超流，而非简单的 $n_0\to0$。
